\documentclass[12pt,a4paper]{ctexart}

\usepackage[a4paper,margin=1.8cm]{geometry}
\usepackage{amsmath,amssymb}
\usepackage{enumitem}
\usepackage{xcolor}
\usepackage{listings}
\usepackage{tcolorbox}
\usepackage{titlesec}
\usepackage{setspace}
\usepackage{hyperref}
% 水印部分
% 直接挂载到 LaTeX 的前景(foreground)渲染层
\AddToHook{shipout/foreground}{
    \begin{tikzpicture}[remember picture, overlay]
        % 在页面正中心放置置顶水印
        \node[text=blue, opacity=0.2, scale=5, rotate=30] 
        at (current page.center) {fifine专升本};
    \end{tikzpicture}
}
\setstretch{1.2}
\setlist[enumerate]{itemsep=0.6em, topsep=0.4em}

\titleformat{\section}{\Large\bfseries}{}{0em}{}
\titleformat{\subsection}{\large\bfseries}{}{0em}{}

\lstset{
    language=Java,
    basicstyle=\ttfamily\small,
    keywordstyle=\color{blue},
    commentstyle=\color{gray},
    stringstyle=\color{red!70!black},
    numbers=left,
    numberstyle=\tiny\color{gray},
    stepnumber=1,
    numbersep=8pt,
    frame=single,
    breaklines=true,
    columns=flexible,
    showstringspaces=false,
    tabsize=4
}

\newtcolorbox{coverbox}{
    colback=gray!5,
    colframe=black!60,
    boxrule=0.8pt,
    arc=2mm,
    left=2mm,right=2mm,top=2mm,bottom=2mm
}

\begin{document}

\begin{center}
    {\LARGE \textbf{专升本 Java 小红书发题题库（第一阶段）}}\\[0.6em]
    {\large 适用方向：专升本学生｜图文发题｜仅题目版｜可用于早题晚解析}\\[0.8em]
    {\normalsize 使用方式：每一节可单独作为 1 篇小红书笔记发布}
\end{center}

\vspace{1em}

\section*{使用说明}
本题库为“小红书发题版”，每组题控制在 3～4 题，便于做成图文轮播。\\
本文件\textbf{不含答案、不含解析}，你可后续单独制作“晚间解析版”。

\newpage

%========================
\section{笔记 1：Java 基础概念查漏补缺}
\begin{coverbox}
\textbf{封面标题建议：}\\
专升本Java｜基础概念这4题，能看出你底子稳不稳\\
\textbf{副标题建议：}\\
别小看基础题，很多同学就丢在这
\end{coverbox}

\subsection*{题目 1}
在 Java 中，以下哪个关键字用于定义常量？
\begin{enumerate}[label=\Alph*.]
    \item \texttt{final}
    \item \texttt{static}
    \item \texttt{abstract}
    \item \texttt{interface}
\end{enumerate}

\subsection*{题目 2}
以下哪个是 Java 中的合法标识符？
\begin{enumerate}[label=\Alph*.]
    \item \texttt{123abc}
    \item \texttt{@abc}
    \item \texttt{\_abc}
    \item \texttt{class}
\end{enumerate}

\subsection*{题目 3}
Java 中用于声明整数类型变量的关键字是
\begin{enumerate}[label=\Alph*.]
    \item \texttt{int}
    \item \texttt{float}
    \item \texttt{double}
    \item \texttt{boolean}
\end{enumerate}

\subsection*{题目 4}
以下哪个是 Java 中的逻辑与运算符？
\begin{enumerate}[label=\Alph*.]
    \item \texttt{\&\&}
    \item \texttt{||}
    \item \texttt{!}
    \item \texttt{\&}
\end{enumerate}

\newpage

%========================
\section{笔记 2：整数运算最容易错的地方}
\begin{coverbox}
\textbf{封面标题建议：}\\
专升本Java｜整数除法这3题，很多人第一眼就错\\
\textbf{副标题建议：}\\
先别急着算，先看数据类型
\end{coverbox}

\subsection*{题目 1}
以下代码的输出结果是
\begin{lstlisting}
int a = 5;
int b = 2;
System.out.println(a / b);
\end{lstlisting}

\begin{enumerate}[label=\Alph*.]
    \item 2.5
    \item 2
    \item 3
    \item 2.0
\end{enumerate}

\subsection*{题目 2}
以下代码的输出结果是
\begin{lstlisting}
int a = 7;
int b = 3;
System.out.println(a % b);
\end{lstlisting}

\begin{enumerate}[label=\Alph*.]
    \item 1
    \item 2
    \item 3
    \item 0
\end{enumerate}

\subsection*{题目 3}
以下代码的输出结果是
\begin{lstlisting}
int a = 8;
int b = 3;
System.out.println(a / b + a % b);
\end{lstlisting}

\begin{enumerate}[label=\Alph*.]
    \item 2
    \item 3
    \item 4
    \item 5
\end{enumerate}

\newpage

%========================
\section{笔记 3：类和对象第一组母题}
\begin{coverbox}
\textbf{封面标题建议：}\\
专升本Java｜类和对象这3题，别只会背概念\\
\textbf{副标题建议：}\\
看懂这组，后面面向对象才不容易乱
\end{coverbox}

\subsection*{题目 1}
以下关于类和对象的说法正确的是
\begin{enumerate}[label=\Alph*.]
    \item 类是对象的实例
    \item 对象是类的模板
    \item 一个类只能创建一个对象
    \item 对象是类的具体实例
\end{enumerate}

\subsection*{题目 2}
以下哪个说法更符合 Java 中对象作为参数传递的特点？
\begin{enumerate}[label=\Alph*.]
    \item 传递的是对象本身
    \item 传递的是对象的引用
    \item 传递的是对象成员变量的值
    \item 传递的是对象类型
\end{enumerate}

\subsection*{题目 3}
下列说法中正确的是
\begin{enumerate}[label=\Alph*.]
    \item 一个对象可以脱离类单独存在
    \item 类用于描述一组对象的共同特征和行为
    \item 对象和类在 Java 中没有区别
    \item 对象只能包含属性，不能包含行为
\end{enumerate}

\newpage

%========================
\section{笔记 4：继承、抽象类、接口辨析题}
\begin{coverbox}
\textbf{封面标题建议：}\\
专升本Java｜extends、abstract、interface 总混？这4题测一下\\
\textbf{副标题建议：}\\
这组题特别适合做一轮查漏补缺
\end{coverbox}

\subsection*{题目 1}
在 Java 中，以下哪个关键字用于实现继承？
\begin{enumerate}[label=\Alph*.]
    \item \texttt{extends}
    \item \texttt{implements}
    \item \texttt{interface}
    \item \texttt{abstract}
\end{enumerate}

\subsection*{题目 2}
以下关于抽象类的说法错误的是
\begin{enumerate}[label=\Alph*.]
    \item 抽象类可以有抽象方法
    \item 抽象类不能被实例化
    \item 抽象类的子类必须实现所有抽象方法
    \item 抽象类中可以有非抽象方法
\end{enumerate}

\subsection*{题目 3}
以下哪个关键字用于定义接口？
\begin{enumerate}[label=\Alph*.]
    \item \texttt{interface}
    \item \texttt{abstract class}
    \item \texttt{class}
    \item \texttt{extends}
\end{enumerate}

\subsection*{题目 4}
以下关于接口的说法错误的是
\begin{enumerate}[label=\Alph*.]
    \item 接口中的方法默认是 \texttt{public abstract} 修饰的
    \item 接口中的变量默认是 \texttt{public static final} 修饰的
    \item 一个类可以实现多个接口
    \item 接口可以继承类
\end{enumerate}

\newpage

%========================
\section{笔记 5：字符串 API 高频题}
\begin{coverbox}
\textbf{封面标题建议：}\\
专升本Java｜String 这4题别再凭感觉做了\\
\textbf{副标题建议：}\\
方法名会背，不等于题能做对
\end{coverbox}

\subsection*{题目 1}
以下哪个方法可以将字符串转换为整数？
\begin{enumerate}[label=\Alph*.]
    \item \texttt{parseInt()}
    \item \texttt{valueOf()}
    \item \texttt{toString()}
    \item \texttt{toInteger()}
\end{enumerate}

\subsection*{题目 2}
在 Java 中，以下哪个方法用于获取字符串的长度？
\begin{enumerate}[label=\Alph*.]
    \item \texttt{length()}
    \item \texttt{size()}
    \item \texttt{count()}
    \item \texttt{getLength()}
\end{enumerate}

\subsection*{题目 3}
以下哪个方法用于将字符串转换为大写？
\begin{enumerate}[label=\Alph*.]
    \item \texttt{toUpperCase()}
    \item \texttt{upperCase()}
    \item \texttt{makeUpperCase()}
    \item \texttt{convertToUpperCase()}
\end{enumerate}

\subsection*{题目 4}
以下哪个方法用于截取字符串？
\begin{enumerate}[label=\Alph*.]
    \item \texttt{substring()}
    \item \texttt{cut()}
    \item \texttt{slice()}
    \item \texttt{truncate()}
\end{enumerate}

\newpage

%========================
\section{笔记 6：异常处理核心考点}
\begin{coverbox}
\textbf{封面标题建议：}\\
专升本Java｜异常处理最容易考的4题，先做一遍\\
\textbf{副标题建议：}\\
finally、catch、常见异常别混了
\end{coverbox}

\subsection*{题目 1}
以下哪个异常类用于处理数组下标越界异常？
\begin{enumerate}[label=\Alph*.]
    \item \texttt{ArrayIndexOutOfBoundsException}
    \item \texttt{NullPointerException}
    \item \texttt{ArithmeticException}
    \item \texttt{ClassCastException}
\end{enumerate}

\subsection*{题目 2}
在 Java 中，以下哪个关键字用于捕获异常？
\begin{enumerate}[label=\Alph*.]
    \item \texttt{try}
    \item \texttt{catch}
    \item \texttt{throw}
    \item \texttt{throws}
\end{enumerate}

\subsection*{题目 3}
以下关于 Java 中异常处理中 \texttt{finally} 块的说法正确的是
\begin{enumerate}[label=\Alph*.]
    \item 无论是否发生异常都会执行
    \item 只有发生异常才会执行
    \item 只有不发生异常才会执行
    \item 只有在 \texttt{try} 块中有 \texttt{return} 语句时才会执行
\end{enumerate}

\subsection*{题目 4}
以下哪种情况更容易导致 \texttt{OutOfMemoryError}？
\begin{enumerate}[label=\Alph*.]
    \item 数组越界
    \item 除数为 0
    \item 内存长期无法释放
    \item 空指针异常
\end{enumerate}

\newpage

%========================
\section{笔记 7：集合框架第一组母题}
\begin{coverbox}
\textbf{封面标题建议：}\\
专升本Java｜List、Set、Map 这4题你能全对吗\\
\textbf{副标题建议：}\\
集合题最怕概念一知半解
\end{coverbox}

\subsection*{题目 1}
以下关于 Java 集合框架的说法正确的是
\begin{enumerate}[label=\Alph*.]
    \item List 允许重复元素
    \item Set 允许重复元素
    \item Map 存储键值对，键可以重复
    \item 以上都不对
\end{enumerate}

\subsection*{题目 2}
在 Java 中，以下哪个集合实现了链表结构？
\begin{enumerate}[label=\Alph*.]
    \item \texttt{ArrayList}
    \item \texttt{LinkedList}
    \item \texttt{HashSet}
    \item \texttt{HashMap}
\end{enumerate}

\subsection*{题目 3}
以下哪个集合是无序且不允许重复元素的？
\begin{enumerate}[label=\Alph*.]
    \item \texttt{ArrayList}
    \item \texttt{HashSet}
    \item \texttt{LinkedHashSet}
    \item \texttt{TreeSet}
\end{enumerate}

\subsection*{题目 4}
以下关于 \texttt{HashMap} 的说法错误的是
\begin{enumerate}[label=\Alph*.]
    \item 存储键值对
    \item 键不允许重复
    \item 是线程安全的
    \item 基于哈希表实现
\end{enumerate}

\newpage

%========================
\section{笔记 8：泛型高频辨析题}
\begin{coverbox}
\textbf{封面标题建议：}\\
专升本Java｜泛型别只会背“提高安全性”，这3题试试\\
\textbf{副标题建议：}\\
一到擦除和边界就开始乱了
\end{coverbox}

\subsection*{题目 1}
以下关于 Java 中泛型的说法错误的是
\begin{enumerate}[label=\Alph*.]
    \item 可以提高代码的安全性
    \item 可以减少类型转换
    \item 泛型类型在运行时会被擦除
    \item 泛型可以直接用于基本数据类型
\end{enumerate}

\subsection*{题目 2}
以下哪个是正确的 Java 泛型声明？
\begin{enumerate}[label=\Alph*.]
    \item \texttt{class MyClass<T extends Number>}
    \item \texttt{class MyClass<T super Number>}
    \item \texttt{class MyClass<T implements Number>}
    \item \texttt{class MyClass<T>}
\end{enumerate}

\subsection*{题目 3}
关于泛型和包装类的说法正确的是
\begin{enumerate}[label=\Alph*.]
    \item \texttt{List<int>} 是合法写法
    \item 泛型中通常应使用包装类而不是基本类型
    \item Java 泛型只支持字符串类型
    \item 泛型和集合无关
\end{enumerate}

\newpage

%========================
\section{笔记 9：反射机制第一轮筛错题}
\begin{coverbox}
\textbf{封面标题建议：}\\
专升本Java｜反射机制这3题，先别被概念绕进去\\
\textbf{副标题建议：}\\
这部分最适合做“筛错题”
\end{coverbox}

\subsection*{题目 1}
以下关于 Java 中反射机制的说法错误的是
\begin{enumerate}[label=\Alph*.]
    \item 可以在运行时获取类的信息
    \item 可以动态创建对象
    \item 不能调用对象的私有方法
    \item 反射机制会降低程序性能
\end{enumerate}

\subsection*{题目 2}
以下关于 Java 中 \texttt{ClassLoader} 的说法错误的是
\begin{enumerate}[label=\Alph*.]
    \item 用于加载类文件
    \item 可以自定义 \texttt{ClassLoader}
    \item 只有一个默认的 \texttt{ClassLoader}
    \item 不同的 \texttt{ClassLoader} 可以加载同名类
\end{enumerate}

\subsection*{题目 3}
关于反射机制，下列说法正确的是
\begin{enumerate}[label=\Alph*.]
    \item 只能在编译期获取类信息
    \item 运行时也可以操作类的成员
    \item 反射只能用于接口
    \item 反射和普通代码执行效率完全相同
\end{enumerate}

\newpage

%========================
\section{笔记 10：I/O 流和 File 高频题}
\begin{coverbox}
\textbf{封面标题建议：}\\
专升本Java｜IO 和 File 这4题，很多人做题顺序就错了\\
\textbf{副标题建议：}\\
先分清“读写”和“字节字符”
\end{coverbox}

\subsection*{题目 1}
以下哪个类用于实现文件读取？
\begin{enumerate}[label=\Alph*.]
    \item \texttt{FileReader}
    \item \texttt{FileWriter}
    \item \texttt{BufferedReader}
    \item \texttt{BufferedWriter}
\end{enumerate}

\subsection*{题目 2}
以下哪个类用于实现文件写入？
\begin{enumerate}[label=\Alph*.]
    \item \texttt{FileReader}
    \item \texttt{FileWriter}
    \item \texttt{BufferedReader}
    \item \texttt{BufferedWriter}
\end{enumerate}

\subsection*{题目 3}
以下关于 Java 中流的说法错误的是
\begin{enumerate}[label=\Alph*.]
    \item 分为字节流和字符流
    \item 字节流以字节为单位处理数据
    \item 字符流以字符为单位处理数据
    \item 字节流和字符流可以直接随意互换
\end{enumerate}

\subsection*{题目 4}
以下哪个方法用于创建一个新文件？
\begin{enumerate}[label=\Alph*.]
    \item \texttt{createNewFile()}
    \item \texttt{mkdir()}
    \item \texttt{mkdirs()}
    \item \texttt{exists()}
\end{enumerate}

\newpage

%========================
\section{笔记 11：JDBC 最容易考的核心题}
\begin{coverbox}
\textbf{封面标题建议：}\\
专升本Java｜JDBC 这4题不熟，后面很容易整章混\\
\textbf{副标题建议：}\\
连接、查询、预编译、事务一次测一下
\end{coverbox}

\subsection*{题目 1}
以下关于 Java 中 JDBC 的说法错误的是
\begin{enumerate}[label=\Alph*.]
    \item 用于连接数据库
    \item 不同数据库有不同驱动
    \item 执行 SQL 语句可以通过 \texttt{Statement} 对象
    \item 不需要关闭连接资源
\end{enumerate}

\subsection*{题目 2}
以下哪个方法用于执行查询语句？
\begin{enumerate}[label=\Alph*.]
    \item \texttt{executeQuery()}
    \item \texttt{executeUpdate()}
    \item \texttt{execute()}
    \item \texttt{runQuery()}
\end{enumerate}

\subsection*{题目 3}
以下关于 Java 中 \texttt{PreparedStatement} 的说法正确的是
\begin{enumerate}[label=\Alph*.]
    \item 性能比 \texttt{Statement} 差
    \item 不能防止 SQL 注入
    \item 可以设置参数
    \item 不支持批量操作
\end{enumerate}

\subsection*{题目 4}
以下关于 Java 中事务的说法错误的是
\begin{enumerate}[label=\Alph*.]
    \item 保证数据的一致性
    \item 可以通过 \texttt{commit()} 提交
    \item 可以通过 \texttt{rollback()} 回滚
    \item 事务中的操作一定能成功
\end{enumerate}

\newpage

%========================
\section{笔记 12：多线程高频辨析题}
\begin{coverbox}
\textbf{封面标题建议：}\\
专升本Java｜多线程这4题，很多人学过但还是会混\\
\textbf{副标题建议：}\\
start、run、sleep、同步一次区分
\end{coverbox}

\subsection*{题目 1}
以下关于线程的说法错误的是
\begin{enumerate}[label=\Alph*.]
    \item 多线程可以提高程序执行效率
    \item 线程之间完全共享所有内存
    \item 实现线程的方式有继承 \texttt{Thread} 类和实现 \texttt{Runnable} 接口
    \item 线程是进程的一部分
\end{enumerate}

\subsection*{题目 2}
在 Java 中，以下哪个方法用于启动线程？
\begin{enumerate}[label=\Alph*.]
    \item \texttt{start()}
    \item \texttt{run()}
    \item \texttt{sleep()}
    \item \texttt{yield()}
\end{enumerate}

\subsection*{题目 3}
以下关于同步的说法中，正确的是
\begin{enumerate}[label=\Alph*.]
    \item 只能在普通方法上使用
    \item 可以使用 \texttt{synchronized} 实现同步
    \item 同步一定会提高程序效率
    \item 同步和线程安全无关
\end{enumerate}

\subsection*{题目 4}
以下哪个方法用于暂停当前线程指定的时间？
\begin{enumerate}[label=\Alph*.]
    \item \texttt{sleep()}
    \item \texttt{wait()}
    \item \texttt{yield()}
    \item \texttt{stop()}
\end{enumerate}

\newpage

%========================
\section{笔记 13：网络编程基础题}
\begin{coverbox}
\textbf{封面标题建议：}\\
专升本Java｜Socket 和 URL 这4题，先把基础拿稳\\
\textbf{副标题建议：}\\
网络编程别一上来就被名词吓住
\end{coverbox}

\subsection*{题目 1}
以下关于 Java 中 Socket 编程的说法错误的是
\begin{enumerate}[label=\Alph*.]
    \item 用于网络通信
    \item \texttt{ServerSocket} 用于服务器端
    \item \texttt{Socket} 用于客户端
    \item 不需要处理异常
\end{enumerate}

\subsection*{题目 2}
以下哪个方法更常用于向输出流发送数据？
\begin{enumerate}[label=\Alph*.]
    \item \texttt{send()}
    \item \texttt{write()}
    \item \texttt{output()}
    \item \texttt{transmit()}
\end{enumerate}

\subsection*{题目 3}
以下关于 Java 中 \texttt{URL} 类的说法错误的是
\begin{enumerate}[label=\Alph*.]
    \item 用于表示统一资源定位符
    \item 可以获取网络资源
    \item 不能解析 URL
    \item 可以获取 URL 的各个部分
\end{enumerate}

\subsection*{题目 4}
以下哪个方法用于获取 URL 的协议部分？
\begin{enumerate}[label=\Alph*.]
    \item \texttt{getProtocol()}
    \item \texttt{getScheme()}
    \item \texttt{getType()}
    \item \texttt{getProtocolType()}
\end{enumerate}

\end{document}